\documentclass{article}
\usepackage{amsmath,amssymb,amsthm}
\usepackage{algorithm}
\usepackage[noend]{algpseudocode}
\usepackage{enumitem}
\usepackage{hyperref}
\usepackage{geometry}
\geometry{margin=1in}
\newtheorem{theorem}{Theorem}
\newtheorem{lemma}{Lemma}
\newtheorem{assumption}{Assumption}
\newtheorem{definition}{Definition}
\newcommand{\E}{\mathbb{E}}
\newcommand{\R}{\mathbb{R}}
\newcommand{\inner}[2]{\langle #1,#2\rangle}
\newcommand{\Breg}[2]{D_{h}(#1,#2)}
\newcommand{\grad}{\nabla}
\begin{document}

\section*{Algorithm Name}
\textbf{Bregman Proximal Gradient (Mirror Descent) with Adaptive Stepsizes}  

Given a convex differentiable objective $f:\R^{n}\rightarrow\R$ and a $1$‑strongly convex prox‑function $h:\R^{n}\rightarrow\R$ (w.r.t. a norm $\|\cdot\|$), the iteration is  

\[
x_{k+1}= \arg\min_{x\in\R^{n}}
\Big\{ \inner{\grad f(x_{k})}{x}+ \frac{1}{\eta_{k}}\,\Breg{x}{x_{k}}\Big\},
\qquad k=0,1,\dots
\tag{1}
\]

where $\eta_{k}>0$ is the stepsize (learning rate) at iteration $k$ and  

\[
\Breg{x}{y}=h(x)-h(y)-\inner{\grad h(y)}{x-y}
\]
is the Bregman divergence induced by $h$.

--------------------------------------------------------------------

\section*{Mathematical Formulation}
\begin{itemize}
    \item \textbf{Input:} initial point $x_{0}\in\R^{n}$, stepsize sequence $\{\eta_{k}\}_{k\ge 0}$.
    \item \textbf{Iterate:} for $k=0,1,2,\dots$
    \[
    \boxed{
    x_{k+1}= \arg\min_{x\in\R^{n}}
    \Big\{ \inner{\grad f(x_{k})}{x}+ \frac{1}{\eta_{k}}\,
    \big( h(x)-h(x_{k})-\inner{\grad h(x_{k})}{x-x_{k}}\big) \Big\}
    }.
    \tag{2}
    \]
    \item \textbf{Termination:} stop when $\|x_{k+1}-x_{k}\|\le \varepsilon$ or after a maximum number of iterations $K_{\max }$.
\end{itemize}

The optimality condition of (2) yields the explicit update (when $h$ is Legendre):
\[
\grad h(x_{k+1}) = \grad h(x_{k}) - \eta_{k}\,\grad f(x_{k}). \tag{3}
\]

--------------------------------------------------------------------

\section*{Pseudocode}
\begin{algorithm}[H]
\caption{Bregman Proximal Gradient (Mirror Descent)}\label{alg:mpg}
\begin{algorithmic}[1]
\Require $x_{0}\in\R^{n}$, stepsize schedule $\{\eta_{k}\}_{k\ge0}$, tolerance $\varepsilon$, max iter $K_{\max}$
\State Compute $u_{0}\gets \grad h(x_{0})$
\For{$k=0$ \textbf{to} $K_{\max}-1$}
    \State $g_{k}\gets \grad f(x_{k})$                         \Comment{gradient of $f$ at $x_{k}$}
    \State $u_{k+1}\gets u_{k} - \eta_{k}\,g_{k}$              \Comment{mirror step, Eq.~(3)}
    \State $x_{k+1}\gets (\grad h)^{-1}(u_{k+1})$              \Comment{map back to primal space}
    \If{$\|x_{k+1}-x_{k}\|\le\varepsilon$}
        \State \textbf{break}
    \EndIf
\EndFor
\State \textbf{return} $x_{k+1}$
\end{algorithmic}
\end{algorithm}

--------------------------------------------------------------------

\section*{Dry Run Example}
Minimize $f(w)=(w-3)^{2}$ on $\R$ using the Euclidean prox $h(w)=\tfrac12 w^{2}$ (so $D_{h}(x,y)=\tfrac12 (x-y)^{2}$).  
Choose $w_{0}=0$, constant stepsize $\eta_{k}=0.1$ for $k=0,1,2$.

\begin{center}
\begin{tabular}{c|c|c|c}
$k$ & $\grad f(w_{k})$ & $w_{k+1}= \arg\min\limits_{w}\{\grad f(w_{k})\,w+\tfrac1{\eta_{k}}D_{h}(w,w_{k})\}$ & $f(w_{k+1})$\\\hline
0 & $2(w_{0}-3)=-6$ & $w_{1}= w_{0}-\eta_{0}\,\grad f(w_{0}) =0-0.1(-6)=0.6$ & $(0.6-3)^{2}=5.76$\\
1 & $2(0.6-3)=-4.8$ & $w_{2}=0.6-0.1(-4.8)=1.08$ & $(1.08-3)^{2}=3.6864$\\
2 & $2(1.08-3)=-3.84$ & $w_{3}=1.08-0.1(-3.84)=1.464$ & $(1.464-3)^{2}=2.3650$
\end{tabular}
\end{center}
The objective value decreases from $9$ (at $w_{0}$) to $2.365$ after three iterations, illustrating the descent property of (1).

--------------------------------------------------------------------

\section*{Assumptions}
\begin{assumption}[Convexity and Smoothness of $f$]\label{as:convex}
The objective $f$ is convex and $L$‑smooth with respect to the norm $\|\cdot\|$, i.e.
\[
f(y)\le f(x)+\inner{\grad f(x)}{y-x}+\frac{L}{2}\|y-x\|^{2},
\qquad\forall x,y\in\R^{n}.
\tag{4}
\]
\end{assumption}

\begin{assumption}[Strong Convexity of $h$]\label{as:hsc}
The prox‑function $h$ is $1$‑strongly convex w.r.t. $\|\cdot\|$:
\[
h(y)\ge h(x)+\inner{\grad h(x)}{y-x}+\frac12\|y-x\|^{2},
\qquad\forall x,y\in\R^{n}.
\tag{5}
\]
Consequently, the Bregman divergence satisfies
\[
\Breg{y}{x}\ge \frac12\|y-x\|^{2}.
\tag{6}
\]
\end{assumption}

\begin{assumption}[Stepsize]\label{as:stepsize}
The stepsizes are constant $\eta_{k}\equiv\eta$, with $0<\eta\le \frac{1}{L}$.
\end{assumption}

\begin{assumption}[Existence of a minimizer]\label{as:min}
There exists at least one optimal point $x^{\star}\in\arg\min_{x}f(x)$ and the Bregman distance to the optimum is bounded:
\[
\Breg{x^{\star}}{x_{0}}\le D_{0}<\infty .
\tag{7}
\]
\end{assumption}

--------------------------------------------------------------------

\section*{Main Convergence Theorem}
\begin{theorem}[Convergence rate of the Bregman proximal gradient]\label{thm:conv}
Under Assumptions \ref{as:convex}–\ref{as:min} and with a constant stepsize $\eta\le 1/L$, the iterates generated by \eqref{alg:mpg} satisfy for all $K\ge 1$  
\[
f(\bar{x}_{K})-f(x^{\star})
\;\le\;
\frac{D_{0}}{\eta K},
\qquad\text{where }\;
\bar{x}_{K}:=\frac{1}{K}\sum_{k=0}^{K-1}x_{k}.
\tag{8}
\]
Consequently, $f(\bar{x}_{K})\to f(x^{\star})$ at the rate $O(1/K)$.
\end{theorem}

--------------------------------------------------------------------

\section*{Theoretical Properties}
\begin{enumerate}[label=\arabic*.]
    \item \textbf{Stability.} Because $h$ is $1$‑strongly convex, the Bregman distance provides a quadratic lower bound (6). This ensures that the mirror step (3) cannot move arbitrarily far, yielding a uniform bound on $\|x_{k+1}-x_{k}\|$.
    \item \textbf{Hyperparameter Sensitivity.} The stepsize must obey $\eta\le 1/L$ (Assumption \ref{as:stepsize}). A larger $\eta$ would break the key inequality (13) in the proof, possibly causing divergence. The bound in (8) is decreasing in $\eta$, so the largest admissible $\eta$ (i.e. $\eta=1/L$) gives the tightest rate.
    \item \textbf{Comparison to Euclidean Gradient Descent.} If $h(w)=\tfrac12\|w\|^{2}$, then $\Breg{x}{y}=\tfrac12\|x-y\|^{2}$ and (3) reduces to the standard gradient descent update $x_{k+1}=x_{k}-\eta\grad f(x_{k})$. Hence Theorem \ref{thm:conv} recovers the classic $O(1/K)$ rate for convex smooth functions.
    \item \textbf{Extension to Non‑Euclidean Norms.} By choosing $h$ as the negative entropy (for the simplex) or a Mahalanobis distance, the same proof works unchanged, illustrating the flexibility of the Bregman framework.
\end{enumerate}

--------------------------------------------------------------------

\section*{Preliminaries (Definitions and Lemmas)}
\begin{definition}[Bregman divergence]
For a differentiable, strictly convex $h$, the Bregman divergence is
\[
\Breg{x}{y}=h(x)-h(y)-\inner{\grad h(y)}{x-y}.
\]
\end{definition}

\begin{lemma}[Three‑point identity]\label{lem:three}
For any $x,y,z\in\R^{n}$,
\[
\inner{\grad h(z)-\grad h(y)}{x-z}
= \Breg{x}{y}-\Breg{x}{z}-\Breg{z}{y}.
\tag{9}
\]
\end{lemma}
\begin{proof}
Expanding the definition of $\Breg{\cdot}{\cdot}$ and rearranging yields (9). \hfill $\square$
\end{proof}

--------------------------------------------------------------------

\section*{Main Proof (Step‑by‑Step Derivation)}
We prove Theorem \ref{thm:conv}. All non‑trivial steps are numbered.

\noindent\textbf{Step 1. Optimality condition of the subproblem.}  
From the definition (2) and the first‑order optimality condition,
\[
\inner{\grad f(x_{k})}{x-x_{k+1}}+\frac{1}{\eta}\inner{\grad h(x_{k+1})-\grad h(x_{k})}{x-x_{k+1}}\ge 0,
\qquad\forall x\in\R^{n}.
\tag{10}
\]

\noindent\textbf{Step 2. Choose $x=x^{\star}$.}  
Plug $x^{\star}$ (optimal point) into (10):
\[
\inner{\grad f(x_{k})}{x^{\star}-x_{k+1}}
+\frac{1}{\eta}\inner{\grad h(x_{k+1})-\grad h(x_{k})}{x^{\star}-x_{k+1}}\ge 0.
\tag{11}
\]

\noindent\textbf{Step 3. Rearrange (11) using the three‑point identity.}  
Apply Lemma \ref{lem:three} with $(x,x_{k+1},x_{k})\gets (x^{\star},x_{k+1},x_{k})$:
\[
\inner{\grad h(x_{k+1})-\grad h(x_{k})}{x^{\star}-x_{k+1}}
= \Breg{x^{\star}}{x_{k}}-\Breg{x^{\star}}{x_{k+1}}-\Breg{x_{k+1}}{x_{k}}.
\tag{12}
\]
Insert (12) into (11):
\[
\inner{\grad f(x_{k})}{x^{\star}-x_{k+1}}
+\frac{1}{\eta}\Big(\Breg{x^{\star}}{x_{k}}-\Breg{x^{\star}}{x_{k+1}}-\Breg{x_{k+1}}{x_{k}}\Big)\ge 0.
\tag{13}
\]

\noindent\textbf{Step 4. Use convexity of $f$.}  
Convexity gives
\[
f(x_{k})-f(x^{\star})\le \inner{\grad f(x_{k})}{x_{k}-x^{\star}}.
\tag{14}
\]
Negating (14) and adding to (13) yields
\[
f(x_{k})-f(x^{\star})
\le \frac{1}{\eta}\Big(\Breg{x^{\star}}{x_{k}}-\Breg{x^{\star}}{x_{k+1}}-\Breg{x_{k+1}}{x_{k}}\Big)
+\inner{\grad f(x_{k})}{x_{k}-x_{k+1}} .
\tag{15}
\]

\noindent\textbf{Step 5. Bound the inner product term using smoothness.}  
From $L$‑smoothness (4) with $y=x_{k+1}$ and $x=x_{k}$,
\[
f(x_{k+1})\le f(x_{k})+\inner{\grad f(x_{k})}{x_{k+1}-x_{k}}+\frac{L}{2}\|x_{k+1}-x_{k}\|^{2}.
\tag{16}
\]
Rearranging,
\[
\inner{\grad f(x_{k})}{x_{k}-x_{k+1}}
\le f(x_{k})-f(x_{k+1})+\frac{L}{2}\|x_{k+1}-x_{k}\|^{2}.
\tag{17}
\]

\noindent\textbf{Step 6. Replace the inner product in (15) by (17).}  
Substituting (17) into (15) gives
\[
f(x_{k})-f(x^{\star})
\le \frac{1}{\eta}\Big(\Breg{x^{\star}}{x_{k}}-\Breg{x^{\star}}{x_{k+1}}-\Breg{x_{k+1}}{x_{k}}\Big)
+f(x_{k})-f(x_{k+1})+\frac{L}{2}\|x_{k+1}-x_{k}\|^{2}.
\tag{18}
\]

\noindent\textbf{Step 7. Cancel $f(x_{k})$ on both sides.}  
After cancellation we obtain
\[
f(x_{k+1})-f(x^{\star})
\le \frac{1}{\eta}\Big(\Breg{x^{\star}}{x_{k}}-\Breg{x^{\star}}{x_{k+1}}-\Breg{x_{k+1}}{x_{k}}\Big)
+\frac{L}{2}\|x_{k+1}-x_{k}\|^{2}.
\tag{19}
\]

\noindent\textbf{Step 8. Relate Bregman distance to the norm.}  
From strong convexity of $h$ (5) we have (6):
\[
\Breg{x_{k+1}}{x_{k}}\ge \frac12\|x_{k+1}-x_{k}\|^{2}.
\tag{20}
\]

\noindent\textbf{Step 9. Use $\eta\le 1/L$ to dominate the quadratic term.}  
Multiply (20) by $1/\eta$ and combine with the last term in (19):
\[
\frac{L}{2}\|x_{k+1}-x_{k}\|^{2}
-\frac{1}{\eta}\Breg{x_{k+1}}{x_{k}}
\le
\Big(\frac{L}{2}-\frac{1}{2\eta}\Big)\|x_{k+1}-x_{k}\|^{2}
\le 0,
\tag{21}
\]
where the inequality follows because $\eta\le 1/L\;\Rightarrow\; \frac{1}{2\eta}\ge \frac{L}{2}$.

\noindent\textbf{Step 10. Drop the non‑positive term.}  
Using (21) in (19) we obtain the clean recursion
\[
f(x_{k+1})-f(x^{\star})
\le \frac{1}{\eta}\Big(\Breg{x^{\star}}{x_{k}}-\Breg{x^{\star}}{x_{k+1}}\Big).
\tag{22}
\]

\noindent\textbf{Step 11. Telescope over $k=0,\dots,K-1$.}  
Summing (22) for $k=0$ to $K-1$ yields
\[
\sum_{k=0}^{K-1}\bigl(f(x_{k+1})-f(x^{\star})\bigr)
\le \frac{1}{\eta}\Big(\Breg{x^{\star}}{x_{0}}-\Breg{x^{\star}}{x_{K}}\Big)
\le \frac{D_{0}}{\eta},
\tag{23}
\]
where $D_{0}$ is defined in (7) and we used $\Breg{x^{\star}}{x_{K}}\ge 0$.

\noindent\textbf{Step 12. Relate the sum to the average iterate.}  
Since $f$ is convex,
\[
f(\bar{x}_{K}) \le \frac{1}{K}\sum_{k=1}^{K} f(x_{k}),
\qquad
\bar{x}_{K}:=\frac{1}{K}\sum_{k=0}^{K-1}x_{k}.
\tag{24}
\]
Combine (24) with (23) (note the index shift) to obtain
\[
K\bigl(f(\bar{x}_{K})-f(x^{\star})\bigr)
\le \frac{D_{0}}{\eta},
\]
which is exactly the bound (8). \hfill $\square$

--------------------------------------------------------------------

\section*{Rate Derivation (With Constants)}
From (8) we have the explicit constant:
\[
f(\bar{x}_{K})-f(x^{\star})\le \frac{D_{0}L}{K},
\qquad\text{since the optimal admissible stepsize is }\eta=\frac{1}{L}.
\tag{25}
\]
Thus the method achieves the optimal $O(1/K)$ rate for convex smooth optimization.

--------------------------------------------------------------------

\section*{Special Cases}
\begin{itemize}
    \item \textbf{Euclidean setup.} $h(w)=\tfrac12\|w\|_{2}^{2}$ gives $D_{h}(x,y)=\tfrac12\|x-y\|_{2}^{2}$ and (3) reduces to ordinary gradient descent. The bound (25) matches the classic result.
    \item \textbf{Probability simplex.} Choosing $h(w)=\sum_{i} w_{i}\log w_{i}$ (negative entropy) yields the Kullback–Leibler Bregman divergence. The algorithm becomes the exponentiated gradient method, and the same $O(1/K)$ guarantee holds with $D_{0}=D_{\mathrm{KL}}(x^{\star}\|x_{0})$.
    \item \textbf{Strongly convex $f$.} If, in addition, $f$ is $\mu$‑strongly convex, a similar derivation (using $f(x_{k})-f(x^{\star})\ge \frac{\mu}{2}\|x_{k}-x^{\star}\|^{2}$) leads to a linear rate $O\big((1-\eta\mu)^{K}\big)$.
\end{itemize}

--------------------------------------------------------------------

\section*{Discussion}
The presented analysis shows that the Bregman proximal gradient (mirror descent) inherits the $O(1/K)$ convergence of classical gradient descent while allowing the geometry of the problem to be encoded via the prox‑function $h$. The proof hinges on three elementary ingredients:

\begin{enumerate}
    \item the optimality condition of the Bregman subproblem,
    \item the three‑point identity of Bregman divergences,
    \item the smoothness‑strong‑convexity balance $\eta\le 1/L$.
\end{enumerate}
No stochasticity or additional variance terms appear; thus the derivation is fully deterministic and elementary, yet it is sufficiently general to cover many non‑Euclidean settings used in modern machine learning (e.g., online learning on the simplex, mirror descent for sparse optimization).

\end{document}